\documentclass[10pt,a4paper]{amsart}
\usepackage[margin=19mm]{geometry}
\usepackage{amsmath,amssymb,mathtools}
\usepackage[scr=rsfs,cal=euler]{mathalfa}
\usepackage{xfrac}
\usepackage[unicode,hidelinks,bookmarks=false]{hyperref}
\newcommand{\ie}{i.e.\ }
\newcommand{\cf}{cf.\ }
\newcommand{\resp}{respectively\ }
\newcommand{\Subsection}{\subsection}
\providecommand{\nobreakdash}{\nobreak\hbox{-}}
\setlength{\emergencystretch}{2em}
\multlinegap0pt
\usepackage{multirow}
\usepackage{mathrsfs}
\usepackage{appendix}
\usepackage{textcomp}
\usepackage{manyfoot}
\usepackage{booktabs}
\usepackage{algorithm}
\usepackage{algorithmicx}
\usepackage{algpseudocode}
\usepackage{listings}
\newtheorem{lemma}{Lemma}
\newtheorem{corollary}{Corollary}
\newtheorem{theorem}{Theorem}%  meant for continuous numbers
\newtheorem{proposition}[theorem]{Proposition}% 
\newtheorem{example}{Example}%
\newtheorem{remark}{Remark}%
\newtheorem{definition}{Definition}%
\begin{document}
\title[A Finitely Generated Module Representation of the Bickley-Na]{A Finitely Generated Module Representation of the Bickley-Naylor Functions}
\author{Anthony Ruffa; Bourama Toni}
\date{}
\maketitle
\noindent\emph{Source and licence.} A Finitely Generated Module Representation of the Bickley-Naylor Functions. Version arXiv:2606.26415v1 (CC BY 4.0). This material is distributed under the Creative Commons Attribution 4.0 International licence, http://creativecommons.org/licenses/by/4.0/, as recorded in the arXiv OAI licence metadata for this exact version and shown on the abstract page https://arxiv.org/abs/2606.26415v1. Full rights notice: licenses/arxiv-2606.26415-CC-BY-4.0.txt.\par\smallskip
\noindent\textit{Editorial scope.} Counted unit: the original \texttt{theorem} environment \texttt{Integrals4340} at source lines 175--182 together with its complete proof at lines 184--233; the delivered document keeps source lines 146--234 so that every referenced label and every citation resolves inside it. Native editable source is the paper's own concatenated TeX; the AMS wrapper is editorial formatting only. No publisher class, upstream script or unrelated artwork is redistributed.
\section{A System of Equations}

Our approach is described in the following. First, we use \emph{Mathematica} to evaluate a multidimensional integral that is an extension of the types of integrals that appeared in \cite{Ruffa2024}. We state the result also as a lemma to our main theorem

\begin{lemma}

\begin{equation}
I
=
\int_0^\infty
\int_0^\infty
\int_0^\infty
e^{
-\frac{x^2}{4}
(x^2+y+z)
}\cdot
e^{-(y+z)^{-1}}\:
dx\:dy\:dz
=
\frac{8\sqrt{\pi}}{x^3}\:
K_2(x).
\label{Integrals4340a}
\end{equation}

\end{lemma}


We then develop a $3\times 3$ system of equations involving Bickley-Naylor functions of consecutive orders 2, 3, and 4. We first prove


\begin{theorem}

\begin{equation}
x^2\:Ki_2(x) + 3x\:Ki_3(x) + 3\:Ki_4(x) = x^2\:K_2(x).
\label{Integrals4340}
\end{equation}
    
\end{theorem}

\begin{proof}
Consider the integral identity in \eqref{Integrals4340a}, i.e.,

\begin{equation}
I
=
\int_0^\infty
\int_0^\infty
\int_0^\infty
e^{
-\frac{x^2}{4}
(x^2+y+z)
}\cdot
e^{-(y+z)^{-1}}\:
dx\:dy\:dz
=
\frac{8\sqrt{\pi}}{x^3}\:
K_2(x).
\label{Integrals4340aa}
\end{equation}

By design, this integral is amenable to the multivariate power substitution approach described and extensively illustrated in \cite{ Ruffa2022,Ruffa2024}. Specifically, we apply the substitutions $x=\sqrt{y+z}\:u_1$, $dx=\sqrt{y+z}\: du_1$, $y=z\:u_2$, and $dy=z\: du_2$ and obtain

\begin{equation}
\begin{split}
I = & \int_0^\infty \int_0^\infty \int_0^\infty e^{-\frac{1}{4}\:z\:x^2(1+u_1^2)(1+u_2)}\cdot
e^{-z^{-1}\:(1+u_2)^{-1} } z^{\frac{3}{2}}\: \sqrt{1+u_2}\: dz\:du_2\:du_1\\
= & \frac{8\sqrt{\pi}}{x^5}
\int_0^\infty
e^{-x\sqrt{1+u_1^2}}\cdot
\frac{ 3+x^2(1+u_1^2)+3\:x\:\sqrt{1+u_1^2}}{\sqrt{(1+u_1^2)^5}}\:du_1.
\label{Integrals4340b}
\end{split}
\end{equation}

The expression \eqref{Integrals4340b} clearly suggests a trigonometric substitution. We further apply the substitution $u_1=\tan\theta$ and $du_1=\sec^2\theta\:d\theta$, which leads to 

\begin{equation}
I
=
\frac{8\sqrt{\pi}}{x^5}
\int_0^\frac{\pi}{2}
e^{-x\sec\theta}\:
(x^2\cos\theta+3x\cos^2\theta+3\cos^3\theta)\:
d\theta.
\label{Integrals4340c}
\end{equation}

To conclude, we equate \eqref{Integrals4340aa} and \eqref{Integrals4340c} and simplify. Thus the claim. \eqref{Integrals4340}.
\end{proof}


\begin{thebibliography}{99}
\bibitem[Ruffa2022]{Ruffa2022} Ruffa, A.A., \& Toni, B. \emph{Innovative Integrals and Their Applications I}, Springer (2022).
\bibitem[Ruffa2024]{Ruffa2024} Ruffa, A.A., \& Toni, B. \emph{Innovative Integrals and Their Applications II}, pp. 82, Springer (2024).
\end{thebibliography}
\end{document}
